\chapter{บทสรุปและข้อเสนอแนะ}
\begin{mypara}
    \indent โครงงานระบบสนับสนุนการตัดสินใจสําหรับการวางแผนระบบขนส่งสาธารณะ 
(DeSS-T: Decision Support System for Public Transit Planning)เป็นเว็บแอปพลิเคชันที่
 เพื่อออกแบบและพัฒนาระบบจำลองที่ช่วยวิเคราะห์ประสิทธิภาพการเดินรถบัส 
โดยนำเทคโนโลยีการจำลองเหตุการณ์แบบไม่ต่อเนื่อง (Discrete-Event Simulation: DES) 
มาประยุกต์ใช้เป็นกลไกหลักในการประมวลผล
\end{mypara}


    \indent ระบบมีการประมวลผลข้อมูลนำเข้า (Input Data) เพื่อสร้างเป็น ฟังก์ชันการแจกแจง (Distribution Function)
สำหรับใช้ในการจำลองแบบสุ่มเชิงสถิติ (Stochastic Simulation) ซึ่งช่วยให้การสุ่มพฤติกรรมการมาถึงของ
ผู้โดยสาร มีความใกล้เคียงกับข้อมูลสถิติจากเหตุการณ์จริงมากที่สุด
จากการดำเนินงานสามารถสรุปผลการศึกษา ปัญหาที่พบ และข้อเสนอแนะได้ดังนี้:
\end{mypara}



\section{สรุปผล}
ระบบ DeSS-T สามารถทำหน้าที่เป็นเครื่องมือสนับสนุนการตัดสินใจผ่านการนำเสนอข้อมูลเชิงสถิติ 
เช่น เวลารอคอยเฉลี่ย และอัตราการใช้งานรถ โดยอาศัยเทคนิคการจำลองเหตุการณ์แบบไม่ต่อเนื่อง 
(Discrete-Event Simulation) เป็นกลไกหลักในการประมวลผล ร่วมกับการใช้ ฟังก์ชันการแจกแจง 
(Distribution Function) เพื่อใช้ในการจำลองแบบสุ่มเชิงสถิติ (Stochastic Simulation) 
สำหรับสุ่มพฤติกรรมการมาถึงของผู้โดยสารให้ใกล้เคียงกับข้อมูลสถิติย้อนหลัง อย่างไรก็ตาม 
ระบบยังมีข้อจำกัดที่สำคัญ ดังนี้:

\begin{itemize}
    \item \textbf{ด้านการนำเข้าและคุณภาพของข้อมูล:} ระบบยังไม่สามารถเชื่อมต่อข้อมูลอัตโนมัติ
    จากอุปกรณ์ภายนอกได้ ทำให้ต้องนำเข้าข้อมูลแบบ Manual ผ่านไฟล์ Excel เท่านั้น ซึ่งอาจทำให้เกิด
    ข้อผิดพลาดในการกรอกข้อมูลและเพิ่มภาระงานของผู้ใช้
    \item \textbf{ข้อจำกัดด้านปริมาณข้อมูล:} จำเป็นต้องมีปริมาณข้อมูลนำเข้าที่มากพอเพื่อ
    ให้ได้ผลลัพธ์ที่เหมาะสม เนื่องจากขนาดกลุ่มตัวอย่าง (Sample Size) ส่งผลโดยตรงต่อความถูกต้อง 
    (Accuracy) และความเที่ยง (Precision) ของระบบ
    \item \textbf{ด้านตัวแปรสภาพแวดล้อม:} ระบบจำลองยังไม่ได้นำปัจจัยภายนอกที่ผันแปรได้ 
    เช่น สภาพภูมิอากาศ หรือความหนาแน่นของการจราจรจริงมาคำนวณร่วมด้วย 
    ซึ่งอาจส่งผลต่อความแม่นยำในสถานการณ์จริง
    \item \textbf{ด้านประสบการณ์ผู้ใช้งาน (User Experience):} ระบบยังมีความซับซ้อนในการใช้งาน
    ระดับหนึ่ง ทำให้ผู้ใช้ต้องใช้เวลาในการเรียนรู้และทำความเข้าใจเครื่องมือก่อนเริ่มใช้งานจริงได้อย่างมีประสิทธิภาพ
    \item \textbf{ด้านการจัดการข้อมูลพื้นฐาน:} ในการตั้งค่าข้อมูลพื้นฐาน (Configuration Data) 
    เมื่อบันทึกแล้วจะไม่สามารถแก้ไขรายละเอียดภายในได้ หากต้องการปรับปรุง ผู้ใช้จำเป็นต้องสร้างชุดข้อมูลใหม่
    ทั้งหมด ซึ่งอาจทำให้เกิดความไม่สะดวกและเพิ่มภาระงานในการจัดการข้อมูลของผู้ใช้
\end{itemize}

\section{ปัญหาที่พบและแนวทางการแก้ไข}
ในการดำเนินงานโครงงานนี้ ได้พบปัญหาหลักหลายประการที่ส่งผลต่อการพัฒนาและการใช้งานระบบ 
ซึ่งได้สรุปปัญหาและแนวทางแก้ไขไว้ดังนี้:

\begin{itemize}
    \item \textbf{ปัญหาด้านการออกแบบระบบ (System Design):}
    \begin{itemize}
        \item \textbf{ปัญหา:} โครงสร้างตัวแปรและโมเดลในช่วงต้นขาดความชัดเจน 
        ทำให้เกิดความไม่สอดคล้องกันระหว่างส่วนประกอบต่าง ๆ เมื่อมีการแก้ไข
        \item \textbf{แนวทางแก้ไข:} ควรลงรายละเอียดโมเดล ให้ครบถ้วนตั้งแต่แรก 
        และมีกาารกำหนดมาตรฐานการตั้งชื่อตัวแปรเพื่อใช้ร่วมกันอย่างเป็นระบบ
    \end{itemize}
    \item \textbf{ปัญหาด้านการบริหารจัดการโครงการ (Project Management):}
    \begin{itemize}
        \item \textbf{ปัญหา:} ในระหว่างการพัฒนาแบบซึ่งได้มีการแบ่งส่วนการทำงาน
        กันแบบคู่ขนาน (Parallel) การสื่อสารไม่เพียงพอจนเกิดปัญหาในการเชื่อมต่อระบบ 
        และการประมาณเวลาคลาดเคลื่อนเนื่องจากแผนงานขาดความละเอียด
        \item \textbf{แนวทางแก้ไข:} ควรมีกาารแบ่งย่อยในการวางแผนดำเนิงานในส่วนขั้นตอนของระยะการพัฒนา
        ให้ละเอียดขึ้น และวางแผนเผื่อเวลาสำหรับอุปสรรคที่ไม่คาดคิด รวมถึงเพิ่มความถี่ในการประชุมและการสื่อสารระหว่างทีม
        เพื่อให้ทุกคนมีความเข้าใจตรงกัน
    \end{itemize}
    \item \textbf{ปัญหาด้านการเก็บรวบรวมข้อมูล (Data Collection):}
    \begin{itemize}
        \item \textbf{ปัญหา:} การเก็บสถิติผู้โดยสารทำได้ลำบากเนื่องจากป้ายหยุดรถบางแห่งไม่มีกล้องวงจรปิด
        \item \textbf{แนวทางแก้ไข:} ใช้วิธีลงพื้นที่สำรวจและเก็บข้อมูลดิบด้วยตนเองเพื่อให้ได้ชุดข้อมูลที่เพียงพอ
        ต่อการสร้าง Distribution 
    \end{itemize}
\end{itemize}

\section{ข้อเสนอแนะและแนวทางการพัฒนาต่อ}
เพื่อให้ระบบ DeSS-T มีความสมบูรณ์ยิ่งขึ้น คณะผู้จัดทำมีข้อเสนอแนะสำหรับการพัฒนาต่อ ดังนี้:

\begin{itemize}
    \item \textbf{การเพิ่มความยืดหยุ่นในการจัดการข้อมูล (Scenario Data):} พัฒนาให้ระบบสามารถ
    แก้ไขลำดับสถานีในแต่ละเส้นทาง (Route) ได้โดยตรง โดยไม่ต้องเลือกสถานีใหม่ทั้งหมด
    \item \textbf{การปรับปรุงการจัดการ Workspace:} พัฒนาให้ข้อมูลพื้นฐาน (Configuration Data) 
    ที่บันทึกไว้สามารถแก้ไขหรือปรับปรุงค่าได้ทันทีแทนการสร้างใหม่
    \item \textbf{การเพิ่มรายละเอียดการแสดงผล:} เพิ่มการแสดงสถิติในส่วนของผู้โดยสารที่ลงจากรถ
     (Alighting Data) รายสถานีให้ชัดเจนขึ้น และเพิ่มการแสดงผลในรูปแบบกราฟิกเพื่อให้เข้าใจง่ายขึ้น
    \item \textbf{การเชื่อมต่อข้อมูลอัตโนมัติ:} พัฒนาการดึงข้อมูลผ่าน API หรือระบบ GPS 
    ของรถบัสโดยตรงเพื่อลดขั้นตอนการกรอกข้อมูล เพิ่มกรรับข้อมูจากกล้องวงจรปิด เพื่อรับข้อมูลผู้โดยสารที่มารอที่ป้าย
    โดยอัตโนมัติและเพิ่มการเชื่อมต่อกับฐานข้อมูลกลางเพื่อให้สามารถอัปเดตข้อมูลได้แบบเรียลไทม์
    \item \textbf{การขยายระบบ Community:} เพิ่มฟังก์ชันให้ผู้ใช้สามารถเลือกใช้พารามิเตอร์หรือ
    การตั้งค่าที่เคยมีผู้อื่นปรับแต่งไว้แล้วมาเป็นต้นแบบในการจำลองได้สะดวกขึ้น โดยการสร้างระบบ Community 
    ที่ให้ผู้ใช้สามารถแชร์และแลกเปลี่ยนข้อมูลการตั้งค่าระหว่างกันได้สะดวกและมีประสิทธิภาพ
\end{itemize}